\documentclass{article}
\usepackage{arxiv}

\usepackage[utf8]{inputenc}
\usepackage[T1]{fontenc}
\usepackage{hyperref}
\usepackage{url}
\usepackage{booktabs}
\usepackage{amsmath, amssymb}
\usepackage{bm}
\usepackage{natbib}
\usepackage{graphicx}

% The emblem lives in the repository-wide assets directory, two levels up.
\graphicspath{{../../assets/}{../../assets/logos/}{../../assets/biota/}{../../assets/diagrams/}}

\title{The Mathematics and Design of FlameSkimmer}

% arxiv.sty labels every paper "A Preprint" in the running head and under the
% title. This is an internal technical report, not a preprint of a submission.
\renewcommand{\headeright}{Technical Report}
\renewcommand{\undertitle}{Technical Report}

% arxiv.sty's title block sets a 24pt strut above the author names and 0.4in
% below them, which leaves the byline stranded between "Technical Report" and
% the date. This is the same layout with those two gaps halved. Overridden here
% rather than edited in arxiv.sty, which is a vendored file.
\makeatletter
\renewcommand{\@maketitle}{%
  \vbox{%
    \hsize\textwidth
    \linewidth\hsize
    \vskip 0.1in
    \@toptitlebar
    \centering
    {\LARGE\sc \@title\par}
    \@bottomtitlebar
    \textsc{\undertitle}\\
    \vskip 0.1in
    \def\And{%
      \end{tabular}\hfil\linebreak[0]\hfil%
      \begin{tabular}[t]{c}\bf\rule{\z@}{12\p@}\ignorespaces%
    }
    \def\AND{%
      \end{tabular}\hfil\linebreak[4]\hfil%
      \begin{tabular}[t]{c}\bf\rule{\z@}{12\p@}\ignorespaces%
    }
    \begin{tabular}[t]{c}\bf\rule{\z@}{12\p@}\@author\end{tabular}%
  \vskip 0.2in \@minus 0.1in \center{\@date}   \vskip 0.2in
  }
}
\makeatother

\author{
  Daniel Boupha Christensen \\
  Environmental, Health, and Safety, University of California, Berkeley \\
  The Green Initiative Fund \\
}

\begin{document}

% Masthead. Licence at the left, group emblem at the right. This is an internal
% report, so the block carries the licence and the project mark only, and
% deliberately no journal, volume, page range or DOI furniture.
\noindent
\begin{minipage}[c]{0.74\textwidth}
  \footnotesize
  Copyright 2024--2026 Strawberry Creek Monitoring Group contributors.\\
  Licensed under the Apache License, Version 2.0.\\
  \url{https://www.apache.org/licenses/LICENSE-2.0}
\end{minipage}%
\hfill
\begin{minipage}[c]{0.16\textwidth}
  \raggedleft
  \includegraphics[width=\linewidth]{SCMGemblem.png}
\end{minipage}

\maketitle

\begin{abstract}
FlameSkimmer is the third detector we built for Strawberry Creek, and its
purpose is to recognize when it is outside the conditions it has learned.
DuskCrayfish gives a point forecast, and CobbleShoal adds a scale that
describes how noisy a sensor usually is, but neither model can indicate that it
has never seen conditions like the ones in front of it. A storm, a change of
season or a long outage receives the same yes or no decision as an ordinary
afternoon. FlameSkimmer keeps CobbleShoal's encoder unchanged and replaces its
forecast head with a small Bayesian neural network, which learns a
distribution over its weights rather than a single value for each. Every
forecast then carries two variances: the noise the sensor is expected to show,
and the disagreement between plausible versions of the model. Here, we set out
the mathematics of that head, how we trained it, how its two variances feed the
detection channels, and what the first trained checkpoint measured. In brief,
the machinery works, but the uncertainty it reports does not yet respond to
storms, for a reason we now understand and can address. Finally, we describe
how this checkpoint is served, experimentally, in the StrawberryWatch dashboard
and how it scored the labelled events there.
\end{abstract}

\section{Introduction}
\label{sec:introduction}

A forecasting detector flags a reading when it lands somewhere the model did not
expect. This works when the model's expectations are good. It fails when the
creek does something the model has rarely or never seen, because the model
still produces a forecast and a scale around it, and reports both with the same
confidence it would have on a dry Tuesday in August. The resulting score
cannot distinguish an unusual reading from an unusual situation.

Storms are the clearest case on this creek. Rain changes conductivity, stage
and temperature together and quickly, rain-mode behaviour is unlabelled at
three of the four nodes, and most of the corpus is dry. Uber faced the same
problem in traffic forecasting, where the departures that matter most, such as
holidays, are the periods with the least history behind them, and for that
reason built a Bayesian network's uncertainty into anomaly
detection~\citep{zhu2017uber}. On this creek, storms play the role of holidays.

Two kinds of uncertainty are involved~\citep{kendall2017uncertainties}.
\emph{Aleatoric} uncertainty is noise in the data itself, such as the
variation a conductivity probe shows even when nothing is happening, and more
data does not reduce it. \emph{Epistemic} uncertainty is uncertainty about the
model. It is large where training data was sparse and small where it was
plentiful, and more data does reduce it. CobbleShoal already estimates the
first through its heteroscedastic scale head. Before FlameSkimmer, nothing in
the repository estimated the second.

\subsection{The trained configuration}
\label{sec:configuration}

Table~\ref{tab:config} lists the settings of the first trained checkpoint,
which we read from the model definition, the training script and the
calibration artifact written beside the weights. The model is experimental.
These are the settings that the evaluation in Section~\ref{sec:evaluation}
measured, and the served model uses the same checkpoint
(Section~\ref{sec:deployment}).

\begin{table}[t]
  \centering
  \caption{The trained configuration of the first FlameSkimmer checkpoint.
    Values above the rule are shared with CobbleShoal's clock model; values
    below it belong to the Bayesian head and its training.}
  \label{tab:config}
  \begin{tabular}{ll}
    \toprule
    Setting & Value \\
    \midrule
    Sites $S$, canonical variables $V$, nodes $N$ & 5, 5, 17 \\
    Window $T$ & 24 timesteps (6 hours at 15 min sampling) \\
    Context features & 2 (hour sine and cosine) \\
    Hidden dimension $d$ & 64 \\
    Seed & 20260806 \\
    \midrule
    Head layers & $128 \to 32 \to 32 \to 2$, tanh between \\
    Prior & $\mathcal{N}(0, 1)$ on every head weight and bias \\
    Initial $\rho$ & $-5$, so $\sigma \approx 0.0067$ \\
    Supervised training targets $N_{\mathrm{train}}$ & 141{,}534 \\
    KL warmup & linear over the first 10\% of steps (752 steps) \\
    Head samples at prediction $S_{\mathrm{w}}$ & 64 \\
    Trainable parameters & 72{,}637, of which 10{,}500 in the head \\
    Operating false alarm rate $q$ & $10^{-3}$ \\
    Calibrated threshold $z_q$ & 16.271 \\
    \bottomrule
  \end{tabular}
\end{table}

The head is small next to the encoder, by design. Its 5{,}250 weights and
biases each carry a mean and a width, which gives the 10{,}500 parameters, and
all of the model's uncertainty lives in those numbers. The 62{,}137 parameters
of the encoder are ordinary point estimates.

\section{Data and Notation}
\label{sec:data}

\subsection{Corpus and split}
\label{sec:corpus}

FlameSkimmer trains on the same windows that CobbleShoal's clock model trained
on. The corpus runs from 4 June 2025 to 29 April 2026, the span over which
every site in the roster reports. The split is chronological, with validation
starting on 1 February 2026. We chose a chronological split because a detector
is deployed forward in time, and adjacent 15 minute windows overlap by 23 of
their 24 steps, so a shuffled split would leak the training set into
validation. Of 19{,}096 fault-free anchors, 14{,}106 lie wholly on the
training side and 4{,}990 on the validation side, and the run used evenly
spaced subsets of 12{,}000 and 3{,}000. The recorded spans are 4 June to 31
December 2025 for training and 25 February to 28 April 2026 for validation.
The scaler is fitted on training rows only, and we hold offline the two
footbridge sensors that the training side never observed, dissolved oxygen and
floating conductivity.

The corpus, the split, the anchors and the calibration routine are all imported
from CobbleShoal's training script rather than copied. A change there reaches
both models, so two copies cannot drift apart and undermine a comparison
between them.

\begin{table}[t]
  \centering
  \caption{Notation. Symbols shared with the CobbleShoal report keep their
    meaning there.}
  \label{tab:notation}
  \begin{tabular}{lll}
    \toprule
    Symbol & Meaning & Dimension \\
    \midrule
    $\bm{h}$ & encoder output for one node, node state and site context & $2d = 128$ \\
    $a$ & last reading of the node, the anchor & scalar \\
    $w$ & one head weight or bias & scalar \\
    $\mu_{w}, \rho_{w}$ & posterior mean and width parameter of $w$ & scalar \\
    $\sigma_{w}$ & posterior standard deviation, $\log(1 + e^{\rho_{w}})$ & scalar \\
    $s_{p}$ & prior standard deviation & scalar, 1 \\
    $q_{\phi}(w)$ & approximate posterior over head weights & \\
    $\beta$ & KL weight during warmup & scalar, $0 \to 1$ \\
    $\mu_{s}, \sigma_{s}$ & forecast and scale under head sample $s$ & $N$ \\
    $\bar{\mu}$ & predictive mean & $N$ \\
    $v_{\mathrm{al}}, v_{\mathrm{ep}}$ & aleatoric and epistemic variance & $N$ \\
    $v_{\mathrm{tot}}$ & total predictive variance & $N$ \\
    $\Lambda$ & combined Fisher statistic & $N$ \\
    \bottomrule
  \end{tabular}
\end{table}

\section{Model Architecture}
\label{sec:architecture}

\subsection{The shared trunk}
\label{sec:trunk}

The encoder is CobbleShoal's NestedEncoder, unchanged: a decay long short-term
memory (LSTM) cell over each sensor node, a sibling exchange within each site,
directed transport downstream, and a final mix back into the nodes. We set out
its mathematics in the CobbleShoal report and do not repeat it here.

For FlameSkimmer, the order of construction matters. We set the global random
seed, construct the encoder, and then set the seed again with a fixed offset
of 10{,}000 before constructing the head. CobbleShoal does the same. At the
same seed the two encoders therefore start from byte-identical weights, which a
unit test checks, so any difference between the two models' results belongs to
their heads rather than to a more favourable draw of the trunk. FlameSkimmer
composes the encoder rather than inheriting CobbleShoal's base class, so none
of the leave-one-out or reconstruction machinery comes with it.

\subsection{The Bayesian head}
\label{sec:head}

The head reads the encoder output $\bm{h}$ for each node and produces two
numbers, a forecast delta and a raw scale. It has three layers,
\begin{equation}
  \bm{z}_{1} = \tanh(W_{1}\bm{h} + \bm{b}_{1}), \qquad
  \bm{z}_{2} = \tanh(W_{2}\bm{z}_{1} + \bm{b}_{2}), \qquad
  (\Delta, r) = W_{3}\bm{z}_{2} + \bm{b}_{3},
\end{equation}
and every entry of every $W$ and $\bm{b}$ is a random variable rather than a
fixed number. The approximate posterior is mean-field Gaussian, so each weight
is independent of the others,
\begin{equation}
  q_{\phi}(w) = \mathcal{N}\big(\mu_{w},\ \sigma_{w}^{2}\big), \qquad
  \sigma_{w} = \log\big(1 + e^{\rho_{w}}\big),
\end{equation}
and the trainable parameters $\phi$ are the pairs $(\mu_{w}, \rho_{w})$. The
softplus keeps $\sigma_{w}$ positive without a constraint. This is the
parameterisation of Bayes by Backprop~\citep{blundell2015weight} and of the
bayesian-torch layers on which we modelled this
head~\citep{krishnan2022bayesiantorch}. The means start where an ordinary
linear layer's weights start, uniform within $\pm 1/\sqrt{\text{fan in}}$, so
the head begins at the same place as CobbleShoal's scale head. Every
$\rho_{w}$ starts at $-5$.

The forecast and its scale are computed as in CobbleShoal,
\begin{equation}
  \mu = a + \Delta, \qquad \sigma = \operatorname{softplus}(r) + 10^{-2},
\end{equation}
with the forecast anchored on the node's last reading and the scale floored so
that the likelihood cannot reward claiming perfect certainty.

\subsection{A Bayesian head on a deterministic trunk}
\label{sec:partial}

Making the encoder Bayesian as well would multiply the number of parameters
that need a width, slow training, and require running the whole encoder once
per sample at prediction time. The evidence suggests it would add little.
Across four inference methods and eight datasets, networks with only some
layers stochastic matched fully stochastic ones, with no systematic benefit
from going further~\citep{sharma2023stochastic}. With a deterministic trunk,
the encoder runs once per window and only the small head is sampled, so 64
samples cost 64 passes through about five thousand weights rather than 64
passes through the whole model.

\subsection{Depth of the head}
\label{sec:depth}

A single hidden layer would be cheaper, but it would fail in a way that matters
here. Mean-field networks with one hidden layer provably cannot raise their
uncertainty in the gap between two well-separated regions of
data~\citep{foong2020expressiveness}. Dry season and storm are, loosely, two
such regions, and the purpose of the model is the uncertainty between and
beyond them. The same work reports similar problems empirically in deeper
networks, which is why we propose a deep ensemble reference in
Section~\ref{sec:limitations}.

\subsection{The heteroscedastic output}
\label{sec:hetero}

The head still predicts a scale for every reading. Replacing the per-reading
scale with weight uncertainty alone would lose the aleatoric part, which the
dispersion channel (Section~\ref{sec:channels}) requires. Combining an
input-dependent noise scale with weight uncertainty in one model follows the
construction of~\citet{kendall2017uncertainties}. For the same reason, we set
aside two off-the-shelf alternatives, post hoc Laplace approximation and
variational Bayesian last layers. As implemented in their libraries, both
learn a single noise level shared by every input, which discards the
per-reading scale.

\section{Training}
\label{sec:training}

\subsection{The evidence lower bound}
\label{sec:elbo}

Training fits $\phi$ by maximising the evidence lower bound (ELBO),
\begin{equation}
  \mathcal{L}_{\mathrm{ELBO}}(\phi) =
    \mathbb{E}_{q_{\phi}(w)}\big[\log p(\mathcal{D} \mid w)\big]
    - \mathrm{KL}\big(q_{\phi}(w)\ \|\ p(w)\big),
\end{equation}
which balances fitting the data against staying close to the prior. The first
term rewards weights that explain the readings. The second term charges for
every bit of certainty the posterior claims beyond what the prior allows. A
weight that the data says nothing about therefore drifts back toward
$\mathcal{N}(0, s_p^2)$ and its width grows, which is the source of epistemic
uncertainty.

The likelihood is CobbleShoal's masked Gaussian negative log likelihood (NLL),
\begin{equation}
  \mathcal{L}_{\mathrm{nll}} = \frac{1}{|M|} \sum_{i \in M}
    \left[ \tfrac{1}{2} \left( \frac{y_i - \mu_i}{\sigma_i} \right)^{2} + \log \sigma_i \right],
\end{equation}
averaged over the entries $M$ where a target was observed and the anchor was
fresh enough to predict a delta from.

\subsection{The summed KL term}
\label{sec:kl}

With a Gaussian posterior and a Gaussian prior, the Kullback-Leibler (KL)
divergence has a closed form per weight,
\begin{equation}
  \mathrm{KL}\big(\mathcal{N}(\mu_{w}, \sigma_{w}^{2})\ \|\ \mathcal{N}(0, s_{p}^{2})\big)
  = \log\frac{s_{p}}{\sigma_{w}} + \frac{\sigma_{w}^{2} + \mu_{w}^{2}}{2 s_{p}^{2}} - \frac{1}{2},
\end{equation}
and the KL of the whole head is the sum of this term over all 5{,}250 weights
and biases. Summing matters. The bayesian-torch base layer returns the mean
over weights instead, which shrinks the penalty by a factor of the parameter
count and turns the ELBO into a different objective without any warning. The
same library also fills its prior width from an argument named
\texttt{prior\_variance}, so a quantity it calls a variance is used as a
standard deviation. FlameSkimmer uses its own small layer rather than depending
on that library, which avoids both problems and also avoids one more import
paying the network file system cost on the OCF host. One unit test checks the
summed KL against PyTorch's own \texttt{kl\_divergence}, and another checks
that it is zero when the posterior equals the prior.

\subsection{Scaling and warmup}
\label{sec:beta}

The likelihood above is a mean over targets, so the KL has to be placed on the
same per-target footing. The loss for one minibatch is
\begin{equation}
  \mathcal{L} = \mathcal{L}_{\mathrm{nll}} + \beta \, \frac{\mathrm{KL}(q_{\phi} \,\|\, p)}{N_{\mathrm{train}}},
\end{equation}
where $N_{\mathrm{train}} = 141{,}534$ is the number of supervised
$(\text{anchor}, \text{node})$ targets in the training set, counted with the
same mask the likelihood uses. Every minibatch therefore carries the same
share of the KL. Blundell et al.\ also propose a schedule that weights the KL
by $2^{M-i}/(2^{M}-1)$ for batch $i$ of $M$, which the BLiTZ library
implements~\citep{esposito2020blitz}. We do not use it, because that schedule
places nearly all of the KL on the first batch of each epoch and almost none
on the rest.

The weight $\beta$ rises linearly from 0 to 1 over the first 10\% of optimiser
steps and then stays at 1. Starting at zero lets the head learn a sensible
forecast before the prior begins to act on it, and holding $\beta$ at 1
afterwards means the trained model optimises the true ELBO, without tempering.

\subsection{Local reparameterization}
\label{sec:localrep}

We estimate the expectation in the ELBO by sampling. The direct approach draws
one weight matrix and passes the whole minibatch through it, but then every
window in the batch shares the same noise, and the gradient estimate does not
improve as the batch grows. The local reparameterization
trick~\citep{kingma2015variational} samples the layer's output instead. For an
input $\bm{x}$ to a layer with independent Gaussian weights, each
pre-activation is itself Gaussian,
\begin{equation}
  b_{j} \sim \mathcal{N}\!\left(\sum_{i} x_{i}\,\mu_{ji} + \mu_{c_j},\ \
    \sum_{i} x_{i}^{2}\,\sigma_{ji}^{2} + \sigma_{c_j}^{2}\right),
\end{equation}
so one standard normal draw per output element gives an exact sample, with
independent noise for every window. The gradient variance then decreases with
batch size, as expected. The training forward pass uses this trick at every
layer.

\subsection{Optimiser and stopping}
\label{sec:optim}

We trained with Adam at a learning rate of $10^{-3}$, batches of 64 anchors and
gradient clipping at 1.0, using CobbleShoal's stopping rule: the learning rate
halves after two epochs without improvement, and training stops after six,
keeping the best epoch. The monitored quantity is the validation negative ELBO
at $\beta = 1$, computed with a fixed noise generator so that the change
between epochs reflects the weights rather than the draw.
Table~\ref{tab:history} shows the run.

\begin{table}[t]
  \centering
  \caption{Training history. NLL is the training likelihood term, KL$/N$ the
    scaled divergence, and the last two columns the mean epistemic share on 600
    training and 600 validation anchors. Epoch 5 was kept.}
  \label{tab:history}
  \begin{tabular}{rrrrrrr}
    \toprule
    Epoch & NLL & KL$/N$ & $\beta$ & Val NLL & Share train & Share val \\
    \midrule
     1 & $-1.699$ & 0.166 & 0.25 & $-1.933$ & 0.160 & 0.147 \\
     2 & $-2.142$ & 0.162 & 0.50 & $-2.110$ & 0.095 & 0.080 \\
     3 & $-2.332$ & 0.159 & 0.75 & $-2.264$ & 0.155 & 0.126 \\
     4 & $-2.215$ & 0.156 & 1.00 & $-2.304$ & 0.131 & 0.122 \\
     5 & $-2.333$ & 0.152 & 1.00 & $-2.315$ & 0.135 & 0.115 \\
     6 & $-2.402$ & 0.148 & 1.00 & $-2.173$ & 0.076 & 0.058 \\
     7 & $-2.431$ & 0.145 & 1.00 & $-2.254$ & 0.123 & 0.091 \\
     8 & $-2.330$ & 0.142 & 1.00 & $-2.182$ & 0.124 & 0.097 \\
     9 & $-2.548$ & 0.140 & 1.00 & $-2.298$ & 0.101 & 0.071 \\
    10 & $-2.652$ & 0.139 & 1.00 & $-2.205$ & 0.129 & 0.099 \\
    11 & $-2.545$ & 0.137 & 1.00 & $-2.182$ & 0.108 & 0.078 \\
    \bottomrule
  \end{tabular}
\end{table}

The first possibility we ruled out was that the KL term dominated the loss, and
it did not come close. Its scaled size of about 0.15 sits beside a likelihood
of about $-2.3$.

\section{Prediction and Detection}
\label{sec:detection}

\subsection{The predictive distribution}
\label{sec:predictive}

At prediction time the encoder runs once and the head runs once for each of
$S_{\mathrm{w}} = 64$ sampled weight sets, giving a forecast $\mu_{s}$ and a
scale $\sigma_{s}$ per sample. The predictive distribution is a mixture of 64
Gaussians, which we summarise by its first two moments,
\begin{align}
  \bar{\mu} &= \frac{1}{S_{\mathrm{w}}} \sum_{s} \mu_{s}, \\
  v_{\mathrm{al}} &= \frac{1}{S_{\mathrm{w}}} \sum_{s} \sigma_{s}^{2}, \qquad
  v_{\mathrm{ep}} = \frac{1}{S_{\mathrm{w}}} \sum_{s} \big(\mu_{s} - \bar{\mu}\big)^{2}, \qquad
  v_{\mathrm{tot}} = v_{\mathrm{al}} + v_{\mathrm{ep}} .
\end{align}
This is the law of total variance. The aleatoric part is the noise each version
of the model expects, averaged over versions, and the epistemic part is how
much the versions disagree about the forecast itself. We use the population
form of the variance, dividing by $S_{\mathrm{w}}$, because the moment-matched
mixture uses that form. The ratio $v_{\mathrm{ep}} / v_{\mathrm{tot}}$, which
we call the \emph{epistemic share}, is the quantity the evaluation tracks.

\subsection{Whole weight samples at prediction time}
\label{sec:weightspace}

Training samples outputs, whereas prediction samples whole weight matrices, and
the difference is intentional. Local reparameterization gives every window in
a batch its own independent noise, which suits training but not scoring: a
window's score would then depend on which other windows happened to share its
batch and where it sat in it. Drawing $S_{\mathrm{w}}$ complete weight sets and
passing every window through the same sets removes that dependence. For any
single window the two schemes give the same distribution, since each layer's
weights are used once per window and are independent of the layers before
them. A unit test checks that a window scores identically alone and stacked
beside another.

\subsection{Fixed serving noise}
\label{sec:crn}

We draw the 64 weight sets once, when the model is constructed, from a fixed
seed and a private generator, and reuse them on every call. Two runs over the
same window therefore give the same score to the last bit. This matters
because an alert that appears on one run and disappears on the next is worse
than either consistent answer. Drawing from a private generator leaves the
global random state untouched, so constructing the model does not disturb
anything seeded after it. We store the standard normal noise rather than the
weights themselves, so the same draw remains valid however training later
moves $\mu$ and $\rho$.

\subsection{Two channels and their variances}
\label{sec:channels}

FlameSkimmer has the two channels that come with a forecast head at no extra
cost, and the most important rule in the design is which variance each channel
reads.

\textbf{Excursion} is the standardised one-step residual against the
\emph{total} predictive standard deviation,
\begin{equation}
  c_{\mathrm{exc}} = \frac{|y - \bar{\mu}|}{\sqrt{v_{\mathrm{tot}}}} .
\end{equation}
A reading far from the forecast under conditions the model knows well should
count for a lot. The same distance under conditions the model reports it does
not know should count for less, and dividing by the total variance has this
effect.

\textbf{Dispersion} reads the \emph{aleatoric} part only,
\begin{equation}
  c_{\mathrm{disp}} = \max\!\left(0,\ \log \frac{\max(\sqrt{v_{\mathrm{al}}}, \varphi)}{\max(\hat{s}, \varphi)}\right),
\end{equation}
where $\hat{s}$ is the observed spread over the last 8 steps and
$\varphi = 10^{-3}$ is a floor. Dispersion tests whether a sensor is flatter
than its own noise, which is how a frozen probe reveals itself. If the
epistemic variance were included, the expected spread would grow whenever the
model was uncertain, an ordinary sensor would then look too flat by comparison,
and every storm would read as a stuck sensor. Both channels are computed by
CobbleShoal's frozen functions, called with different scales. We do not modify
their arithmetic, because the calibrated artifacts depend on it.

\subsection{Nulls, combination and threshold}
\label{sec:calibration}

Every step after the channels follows CobbleShoal's procedure, run on
FlameSkimmer's scores. We convert each channel to a right-tail probability
against its own fault-free distribution, empirical below the 95th percentile
and a fitted Generalised Pareto distribution above it~\citep{pickands1975gpd}.
We combine the two probabilities per node as
$\Lambda = -2 \sum_{j} \log p_{j}$~\citep{fisher1925methods}, and the
threshold comes from peaks over threshold on the fault-free values of
$\Lambda$~\citep{siffer2017spot},
\begin{equation}
  z_{q} = u - \sigma \log\!\left(\frac{q\,n}{N_{u}}\right) \quad (\gamma = 0).
\end{equation}
The fit used 400 fault-free training windows, giving $n = 4{,}769$ node scores,
an initial threshold $u = 10.135$ with $N_{u} = 96$ exceedances,
$\sigma = 2.044$ and $\gamma = 0$, so $z_{q} = 16.271$ at $q = 10^{-3}$. On
2{,}118 held-out validation scores, the measured exceedance rate was
$4.7 \times 10^{-4}$.

As for CobbleShoal, we handle rain outside the model. The rain gate moves the
threshold and never the score, and FlameSkimmer declares no built-in support
behaviour. The comparison in the next section uses the ungated threshold for
both models.

\section{Evaluation}
\label{sec:evaluation}

We compared FlameSkimmer with CobbleShoal's clock model, trained by the same
script on the same split from the same seed. Every number is in the scaler's
standardised units, and the evaluation checks that both models were scaled
identically before comparing them. We report results per node, because a mean
across nodes hides the nodes where a model is worse.

\subsection{Forecast quality}
\label{sec:quality}

Table~\ref{tab:quality} gives the NLL, the continuous ranked probability score
(CRPS) and the coverage of the central 90\% interval on the 3{,}000 validation
anchors. The CRPS uses the closed form for a Gaussian~\citep{gneiting2007scoring},
applied to the moment-matched predictive distribution, and we copied it from
uncertainty-toolbox~\citep{chung2021toolbox} rather than adding the package.
The NLL here includes the $\tfrac{1}{2}\log 2\pi$ constant that the training
loss leaves out and is taken over the full predictive distribution, so it is
not directly comparable with the validation NLL in Table~\ref{tab:history}.

\begin{table}[t]
  \centering
  \caption{Forecast quality on validation, FlameSkimmer against CobbleShoal.
    Lower NLL and CRPS are better; the 90\% interval should cover close to
    90\%. Footbridge has only 21 supervised targets on the validation side.}
  \label{tab:quality}
  \begin{tabular}{lrrrrrrr}
    \toprule
     & & \multicolumn{2}{c}{NLL} & \multicolumn{2}{c}{CRPS} & \multicolumn{2}{c}{Coverage 90} \\
    Node & $n$ & Flame & Cobble & Flame & Cobble & Flame & Cobble \\
    \midrule
    all nodes                 & 31{,}650 & $-1.460$ & $-1.157$ & 0.0335 & 0.0330 & 0.972 & 0.961 \\
    \midrule
    north\_fork\_0 cond.      & 2{,}996 & $-1.588$ & $-1.410$ & 0.0243 & 0.0277 & 0.972 & 0.977 \\
    north\_fork\_0 depth      & 2{,}996 & $-0.503$ &  $0.461$ & 0.0470 & 0.0522 & 0.977 & 0.988 \\
    north\_fork\_0 temp.      & 2{,}996 & $-1.432$ &  $0.275$ & 0.0178 & 0.0154 & 0.983 & 0.894 \\
    footbridge cond.          &      21 & $-1.222$ &  $0.045$ & 0.0383 & 0.1041 & 1.000 & 1.000 \\
    footbridge depth          &      21 & $-1.052$ & $-0.590$ & 0.0490 & 0.0650 & 0.952 & 1.000 \\
    footbridge temp.          &      21 & $-1.712$ & $-1.443$ & 0.0220 & 0.0252 & 1.000 & 1.000 \\
    south\_fork\_1 cond.      & 1{,}536 & $-1.382$ & $-0.681$ & 0.0330 & 0.0543 & 0.958 & 0.980 \\
    south\_fork\_1 depth      & 1{,}536 & $-0.540$ & $-0.530$ & 0.0701 & 0.0688 & 0.993 & 0.991 \\
    south\_fork\_1 temp.      & 1{,}536 & $-1.645$ & $-1.755$ & 0.0235 & 0.0206 & 0.990 & 0.978 \\
    south\_fork\_2 cond.      & 2{,}997 & $-1.924$ & $-1.664$ & 0.0217 & 0.0258 & 0.973 & 0.939 \\
    south\_fork\_2 depth      & 2{,}997 & $-0.903$ & $-1.030$ & 0.0733 & 0.0496 & 0.989 & 0.985 \\
    south\_fork\_2 temp.      & 2{,}997 & $-2.115$ & $-2.336$ & 0.0155 & 0.0125 & 0.960 & 0.934 \\
    oxford cond.              & 3{,}000 & $-1.980$ & $-1.894$ & 0.0157 & 0.0166 & 0.984 & 0.985 \\
    oxford depth              & 3{,}000 & $-0.972$ & $-0.899$ & 0.0570 & 0.0590 & 0.907 & 0.968 \\
    oxford temp.              & 3{,}000 & $-2.141$ & $-2.184$ & 0.0160 & 0.0148 & 0.985 & 0.943 \\
    \bottomrule
  \end{tabular}
\end{table}

The results are mixed. FlameSkimmer is better on likelihood overall, and better
on both NLL and CRPS at every conductivity node, which is the variable scored
for detection. It is worse on NLL at four nodes and on CRPS at six, mostly
temperature and stage on the south fork and at oxford, so it does not meet the
proposed criterion that no node be worse on either measure. Its intervals are
also too wide. The central 90\% interval covers 97.2\% of validation targets,
outside the 85 to 95\% band that we would count as well calibrated. Wide
intervals improve the likelihood on rough nodes and cost CRPS on smooth ones,
which is consistent with the pattern in the table.

\subsection{Epistemic behaviour}
\label{sec:epistemic}

The test the model exists to pass is whether its epistemic share rises where it
has seen little data. We count a step as rain when more than 0.1 mm fell over
the preceding six hours, according to the Open-Meteo archive.
Table~\ref{tab:epistemic} gives the share on every supervised target, split by
side and by rain, together with the median epistemic standard deviation in
absolute terms.

\begin{table}[t]
  \centering
  \caption{Epistemic share $v_{\mathrm{ep}} / v_{\mathrm{tot}}$ per supervised
    target, and the median epistemic standard deviation
    $\sqrt{v_{\mathrm{ep}}}$. The three event rows count rain steps only; the
    January event is the botanical garden actuator.}
  \label{tab:epistemic}
  \begin{tabular}{lrrrr}
    \toprule
    Group & $n$ & Median share & Quartiles & Median $\sqrt{v_{\mathrm{ep}}}$ \\
    \midrule
    Training, dry      & 131{,}459 & 0.132 & 0.082 to 0.215 & 0.0149 \\
    Training, rain     &  10{,}075 & 0.066 & 0.008 to 0.151 & 0.0147 \\
    Validation, dry    &  29{,}385 & 0.132 & 0.049 to 0.184 & 0.0149 \\
    Validation, rain   &   2{,}265 & 0.027 & 0.002 to 0.117 & 0.0147 \\
    \midrule
    Event, November 2025 rain & 813 & 0.021 & 0.001 to 0.084 & \\
    Event, April 2026 rainfall & 810 & 0.064 & 0.003 to 0.157 & \\
    Event, January 2026 actuator & 5{,}664 & 0.022 & 0.002 to 0.092 & \\
    \bottomrule
  \end{tabular}
\end{table}

The share falls in rain rather than rising, on both sides of the split and in
all three event windows. The last column explains why. The absolute epistemic
standard deviation is 0.0149 or 0.0147 in every group, so the epistemic
variance is in effect a constant. What changes in rain is the aleatoric part:
the scale head correctly predicts a noisier sensor, the total variance grows,
and a constant divided by a larger total becomes smaller. Thus, the model is
not reporting greater certainty in storms. Its epistemic variance carries no
information about conditions at all.

\subsection{Detection}
\label{sec:detectionresults}

Table~\ref{tab:events} scores every labelled event in the catalogue. For each
entry, we compared the peak combined score at the named site's conductivity
node, over every anchor in the event window, with each model's threshold at
$q = 10^{-3}$: 16.271 for FlameSkimmer and 29.070 for CobbleShoal. The two
thresholds are on different scales, since one combines two channels and the
other combines four.

\begin{table}[t]
  \centering
  \caption{Peak score at the graded node in each labelled event window. Entries
    labelled relative only are reported, not graded.}
  \label{tab:events}
  \begin{tabular}{llrrl}
    \toprule
    Event & Site & Flame & Cobble & Label \\
    \midrule
    June 2025 discharge         & south\_fork\_1 & 17.655 & 34.539 & anomaly, both flagged \\
    September 2025 overnight    & south\_fork\_1 & 23.740 & 47.566 & anomaly, both flagged \\
    September 2025 overnight    & south\_fork\_2 & 15.796 & 28.405 & anomaly, both missed \\
    March 2026 hydrant          & north\_fork\_0 & 19.467 & 46.026 & anomaly, both flagged \\
    \midrule
    June 2025 discharge         & south\_fork\_2 & 17.880 & 26.872 & relative only \\
    August 2025 sprinklers      & north\_fork\_0 & 11.228 & 13.507 & relative only \\
    November 2025 foam          & north\_fork\_0 & 23.387 & 48.388 & relative only \\
    \bottomrule
  \end{tabular}
\end{table}

Both models flag three of the four labelled anomalies, and both miss the
south\_fork\_2 half of the September event, FlameSkimmer by about 3\% of its
threshold. FlameSkimmer's margins are thinner, however. Its peaks sit 8 to
46\% above its threshold, whereas CobbleShoal's sit 19 to 64\% above its own,
which is the expected cost of dropping the leave-one-out and reconstruction
channels.

\subsection{False alarm rate}
\label{sec:far}

Table~\ref{tab:far} gives the rate at which fault-free validation node scores
crossed each model's threshold at $q = 10^{-3}$, overall and split by rain.

\begin{table}[t]
  \centering
  \caption{Empirical false alarm rate on fault-free validation anchors at
    $q = 10^{-3}$. CobbleShoal has more finite scores because its leave-one-out
    channel does not need a fresh anchor.}
  \label{tab:far}
  \begin{tabular}{lrrrrrr}
    \toprule
     & \multicolumn{2}{c}{All} & \multicolumn{2}{c}{Dry} & \multicolumn{2}{c}{Rain} \\
    Model & Rate & $n$ & Rate & $n$ & Rate & $n$ \\
    \midrule
    FlameSkimmer & $1.07 \times 10^{-3}$ & 31{,}722 & $7.5 \times 10^{-4}$ & 29{,}445 & $5.3 \times 10^{-3}$ & 2{,}277 \\
    CobbleShoal  & $1.36 \times 10^{-3}$ & 37{,}626 & $6.9 \times 10^{-4}$ & 34{,}905 & $9.9 \times 10^{-3}$ & 2{,}721 \\
    \bottomrule
  \end{tabular}
\end{table}

The overall rate is close to nominal for both models, well within the order of
magnitude that the peaks-over-threshold diagnostics already treat as
acceptable. In rain, FlameSkimmer flags at about half CobbleShoal's rate. This
resembles the benefit the model was built for, but
Section~\ref{sec:epistemic} rules out the intended explanation, since the
epistemic variance does not change in rain. The more likely cause is the wider
total intervals reported in Section~\ref{sec:quality}, which damp the
excursion channel everywhere and matter most where readings move most.

\subsection{The gate criterion}
\label{sec:gate}

The plan set one test of whether the Bayesian head does its job: the median
epistemic share on validation rain steps should be at least 1.5 times the
median on training dry steps. The measured ratio is $0.027 / 0.132 = 0.20$.

The cause lies in the trained widths. Table~\ref{tab:sigma} gives the
posterior standard deviations after training. Every one of them started at
$\log(1 + e^{-5}) \approx 0.0067$, and they have changed little.

\begin{table}[t]
  \centering
  \caption{Posterior standard deviations of the head after training. Every value
    started at 0.0067.}
  \label{tab:sigma}
  \begin{tabular}{lrrr}
    \toprule
    Layer & Median $\sigma$, weights & Largest $\sigma$, weights & Median $\sigma$, biases \\
    \midrule
    $128 \to 32$ & 0.0129 & 0.0160 & 0.0093 \\
    $32 \to 32$  & 0.0064 & 0.0108 & 0.0061 \\
    $32 \to 2$   & 0.0053 & 0.0068 & 0.0036 \\
    \bottomrule
  \end{tabular}
\end{table}

The ceiling on the first layer is not a coincidence. Adam moves each parameter
by at most roughly the learning rate per step, and the kept epoch came after
about 940 steps at a learning rate of $10^{-3}$. Starting from $-5$, no
$\rho_{w}$ could have climbed much past $-4.06$, which corresponds to
$\sigma_{w} \approx 0.017$. The largest width learned is 0.016. Thus, the
optimiser and the starting point limited the widths, rather than the data, and
early stopping kept the epoch immediately after $\beta$ reached 1, before the
prior had any time at full strength.

This also explains why the epistemic variance is constant. When every width
sits near the same small value, the variance the head adds is roughly that
width squared times the squared size of the activations feeding each layer, and
after a tanh those activations have much the same size for almost every input.
The resulting variance depends little on the window, so it cannot rise in a
storm.

These results do not show that a Bayesian head cannot do the job. They show
that this head was not trained long enough, or started too narrow, for its
widths to learn anything.

\section{Deployment}
\label{sec:deployment}

\subsection{Serving}
\label{sec:serving}

Since 16 October 2026, FlameSkimmer can be switched on in the StrawberryWatch
dashboard. It is off by default, and the dashboard shows the notice that the
model is experimental, that only one version has been trained, and that its
uncertainty estimate failed its own test (Section~\ref{sec:gate}). The served
model is the checkpoint described in this report, held to its calibrated
threshold $z_q = 16.271$ at $q = 10^{-3}$, ten times the rate CobbleShoal is
served at.

FlameSkimmer is served through CobbleShoal's detector, with the same combined
Fisher rule, the same two day checks and the same rain gate, which is on by
default. Two further steps reproduce how its calibration was fitted. (1)
Readings outside the plausible ranges recorded with the calibration corpus are
dropped. In each probe's reporting units, these are 0 to 40 for temperature, 0 to 5{,}000 for conductivity and floating
conductivity, $-5{,}000$ to 5{,}000 for depth, and $-5$ to 200 for dissolved
oxygen. (2) The two footbridge sensors that training never observed,
dissolved oxygen and floating conductivity, are held offline as they were in
every calibration window, and their scores are left blank so they can neither
fire nor appear.

\subsection{Labelled events}
\label{sec:servedevents}

On 25 October 2026 we scored the four labelled anomalies with the served model,
at every anchor of each event, with the rain gate on. Table~\ref{tab:served}
gives the peak combined statistic at the labelled site's conductivity node,
beside the peaks in Table~\ref{tab:events}.

\begin{table}[t]
  \centering
  \caption{Peak combined statistic at the labelled site's conductivity node,
    from the evaluation in Table~\ref{tab:events} and from the served model.
    The share is the largest fraction of its rain-gated line any anchor
    reached in serving.}
  \label{tab:served}
  \begin{tabular}{llrrrl}
    \toprule
    Event & Site & Evaluation & Served & Share of line & Flagged \\
    \midrule
    June 2025 discharge      & south\_fork\_1 & 17.655 & 16.099 & 0.989 & no \\
    September 2025 overnight & south\_fork\_1 & 23.740 & 23.740 & 1.437 & yes \\
    September 2025 overnight & south\_fork\_2 & 15.796 & 15.796 & 0.745 & no \\
    March 2026 hydrant       & north\_fork\_0 & 19.467 & 19.467 & 1.196 & yes \\
    \bottomrule
  \end{tabular}
\end{table}

The served model reproduced three of the four evaluation peaks exactly and
flagged two of the four anomalies, against three in the evaluation. The
difference is the June discharge. The dashboard floors every timestamp to its
quarter hour before the model reads it, whereas the training corpus matched
each reading to the nearest grid point within half a step. All 559
south\_fork\_1 readings in the June event file are off the quarter hour, and
under the dashboard's flooring the peak fell from 17.655 to 16.099, just below
the line of 16.271. The same flooring applies to CobbleShoal. In the
September event at south\_fork\_2, the rain gate raised the line, and the
share of line fell from 0.971 without the gate to 0.745 with it.

\section{Limitations and Future Work}
\label{sec:limitations}

The first thing to address is the training of the widths. Three changes target
Section~\ref{sec:gate} directly: (1) start $\rho$ less negative, for example at
$-3$, which gives $\sigma \approx 0.049$; (2) give the $\rho$ parameters their
own, larger learning rate; and (3) stop early only after $\beta$ has been at 1
for several epochs. Each is a change to a default rather than to the design,
and we recommend measuring the gate again after these changes, before anything
else is built on the model.

The plan also calls for a reference when the gate fails: a five-seed deep
ensemble of the same trunk with CobbleShoal's deterministic scale head,
reported beside FlameSkimmer. Ensembles match or outperform approximate
Bayesian networks in most comparisons~\citep{lakshminarayanan2017ensembles}
and held up best under dataset shift in the largest benchmark of the
question~\citep{ovadia2019shift}, although that benchmark is classification.
The ensemble is not Bayesian over weights, but it is the bar a Bayesian head
should clear, and we have not yet run it. Monte Carlo dropout, the other
low-cost option, was the weakest of four methods in a rainfall-runoff LSTM
benchmark~\citep{klotz2022rainfall}, and we do not plan to use it.

We expected the chronological split to put most of the rain on the validation
side. By anchor count it does not: 7.0\% of training anchors and 7.3\% of
validation anchors had rain in the preceding six hours. The storms may still
differ in size or season between the two sides, but the split is a weaker test
of shift than we assumed, and a blocked split or a held-out storm would give a
sharper one.

We have trained only one seed. Every result in Section~\ref{sec:evaluation}
should be repeated across several seeds, and reported per seed, before it is
relied on.

The event catalogue has no true negatives, that is, no storm that the group has
labelled as a known non-event. The rain false alarm comparison therefore rests
on fault-free statistics alone, and one labelled storm would turn it into a
real test case.

FlameSkimmer is served experimentally (Section~\ref{sec:deployment}),
although the evaluation in Section~\ref{sec:evaluation} has not been accepted
and the gate in Section~\ref{sec:gate} failed. The plan made serving
conditional on that evaluation; serving went ahead so that the model could be
tried in the dashboard, and it stays off by default. The README's model table
describes FlameSkimmer, while the user manual has not yet been updated.
Whether the epistemic signal should ever become an alert type of its own is a
separate decision, and nothing in this design adds an uncertainty rule to the
alerting path.

The dashboard's quarter-hour flooring moved one labelled peak below its line
(Section~\ref{sec:servedevents}). We recommend either matching readings to the
grid in serving the way the training corpus did, or recalibrating on floored
readings, so that the served scores are drawn from the distribution the
threshold was fitted on.

Finally, every number in this document describes one experimental checkpoint
and the calibration fitted alongside it. These numbers will move when the model
is retrained, and some of them are intended to.

\bibliographystyle{unsrtnat}
\bibliography{references}

\end{document}
